\section*{الفصل \ref{ch:g6:life-through-seasons} --- الحياة عبر فصول السنة}

\begin{solution}{exo:g6:life-through-seasons:1}
طول النهار (\omterm{def:g6:exploring-our-environment:conditions}{الضوء})، \omterm{def:g6:exploring-our-environment:conditions}{ودرجة الحرارة}، وندى التربة أو الهواء تبعًا
للمطر والجفاف.
\end{solution}

\begin{solution}{exo:g6:life-through-seasons:2}
\omterm{def:g6:life-through-seasons:trees}{النفضية} --- تُسقط أوراقها كلها في الخريف: البلوط، والزان، وكستناء
الحصان. \omterm{def:g6:life-through-seasons:trees}{ودائمة الخضرة} --- تُبقي أوراقًا عاملة طول الشتاء: الصنوبر،
والبهشية.
\end{solution}

\begin{solution}{exo:g6:life-through-seasons:3}
فرخ كامل صغير --- \omterm{def:g1:plants-around-us:parts}{أوراق} صغيرة مطوية ملفوفة في حراشف تصدّ الطقس.
وقد بُني في الصيف السابق.
\end{solution}

\begin{solution}{exo:g6:life-through-seasons:4}
الحولي يموت كله ويعبر الشتاء بذورًا. والنرجس يموت أعلاه بينما
ينتظر \omterm{def:g1:plants-around-us:plant}{النبات} تحت الأرض بصلةً.
\end{solution}

\begin{solution}{exo:g6:life-through-seasons:5}
لأن الأفرخ داخل براعمه كانت قد اكتملت في الصيف الماضي؛ ولم يكن
ينقص \omterm{def:g1:plants-around-us:tree}{الغصن} إلا الدفء، وهو ما يوفّره ماء البيت قبل الحديقة
بأسابيع.
\end{solution}

\begin{solution}{exo:g6:life-through-seasons:6}
\omterm{def:g3:animals-through-seasons:migration}{الهجرة}: السنونو. \omterm{def:g3:animals-through-seasons:hibernation}{والسبات}: القنفذ (أو الحلزون المختوم). والبقاء
نشطًا: الثعلب، والعقعق، وأبو الحنّاء. والشتاء في المراحل الصغيرة ---
بيضًا أو \omterm{ex:g2:animals-grow-up:butterfly}{يرقة} أو شرنقة: أكثر \omterm{prop:g3:sorting-living-things:groups}{الحشرات}، كالجندب أو الفراشة.
\end{solution}

\begin{solution}{exo:g6:life-through-seasons:7}
الفراشة: شرنقة تحت حجر قمة الجدار. والجندب: \omterm{def:g2:animals-grow-up:born}{بيض} في التربة.
والدعسوقة: بالغة نائمة متجمعة في عمق فرش \omterm{def:g1:plants-around-us:parts}{الأوراق}. والحلزون: مختوم
في صدفته في الضفة. والسنونو: في إفريقيا، بذاته المهاجرة.
\end{solution}

\begin{solution}{exo:g6:life-through-seasons:8}
ليأتي أي فرق بين الجداول الأربعة من المواسم لا من المسح --- وهي
قاعدة الاختبار المنصف في
\cref{met:g6:exploring-our-environment:survey}، مطبَّقة عبر الزمن
لا عبر المواقع.
\end{solution}

\begin{solution}{exo:g6:life-through-seasons:9}
أولًا: كثير من ``الغائبين'' حاضر في صور مختفية --- \omterm{def:g2:animals-grow-up:born}{بيض}، وشرانق،
وبالغات نائمة، وأبصال --- يفوّتها مسح يمشي من غير حفر ورفع.
وثانيًا: بعضهم ليس مفقودًا بل في مكان آخر: فالمهاجرون سيعودون.
فقائمة كانون الثاني تبخس المقيمين وتعدّ المسافرين زائلين.
\end{solution}

\begin{solution}{exo:g6:life-through-seasons:10}
\omterm{def:g1:plants-around-us:plant}{النبات} في حال مخزنه تحت الأرض --- حي، منتظر أبصالًا. وفي آذار
سترسل الأبصال أوراقًا وأزهارًا من جديد، أكثف إن تكاثرت الكتل.
\end{solution}

\begin{solution}{exo:g6:life-through-seasons:11}
ذهب إلى أين، وبأي صورة؟ الضفادع تشتّي خامدة --- مدفونة في قاع
البركة الطيني أو في مآوٍ ندية على اليابسة --- تعبر شهور البرد
مخدَّرة حتى يدعوها الربيع إلى الماء للتوالد.
\end{solution}

\begin{solution}{exo:g6:life-through-seasons:12}
مشكلة الشتاء عند \omterm{def:g1:plants-around-us:parts}{الورقة} العاملة هي الاحتفاظ بالماء: فالأرض
المتجمدة لا تعطي \omterm{def:g1:plants-around-us:parts}{الجذور} شيئًا يُذكر، بينما تظل الريح والشمس
تسحبانه --- إنها مشكلة جفاف حشرة القمل الخشبي، الشتاء كله. وشمع
\omterm{def:g1:plants-around-us:parts}{ورقة} \omterm{def:g6:life-through-seasons:trees}{دائمة الخضرة} وبنيتها الإبرية يقلّصان تلك الخسائر إلى ما
تستطيع الشجرة احتماله. أما \omterm{def:g1:plants-around-us:parts}{الورقة} العريضة الرقيقة فلا يمكن ختمها
هكذا؛ وجواب الشجرة \omterm{def:g6:life-through-seasons:trees}{النفضية} الأرخص أن تشطب أوراقها الرقيقة كل خريف
وتخزّن الربيع \omterm{def:g6:life-through-seasons:buds}{براعم} بدل ذلك.
\end{solution}

\begin{solution}{pb:g6:life-through-seasons:1}
\textbf{1.} \omterm{def:g6:life-through-seasons:buds}{البراعم} على كل \omterm{def:g1:plants-around-us:tree}{غصن} --- وهي المخزن الشتوي الذي يحمل
أفرخ الربيع القادم، وقد بُنيت في الصيف السابق.

\textbf{2.} بين صورتي آذار ونيسان (مع طول الأيام واعتدال الصباحات
في الملاحظات). والشجرة تجيب عن ارتفاع الحرارة وطول \omterm{def:g6:exploring-our-environment:conditions}{الضوء}.

\textbf{3.} \omterm{prop:g3:sorting-living-things:groups}{الحشرات} المؤبِّرة --- والنحل قبل كل شيء --- يجب أن
تزور الأزهار، ليبلغ اللقاح المدقات ويتبعه \omterm{def:g5:flower-to-fruit:steps}{الإخصاب}. \omterm{prop:g3:flowers-fruits-seeds:transform}{والثمرة} \omterm{def:g2:from-seed-to-plant:seed}{بذرة}
تكوّنت من بييضة مخصَّبة، في غلافها الشوكي.

\textbf{4.} نفضية. وربيعها القادم يركب صور الشتاء براعمَ على
\omterm{def:g1:plants-around-us:tree}{الأغصان} العارية.

\textbf{5.} تحت العشب، أبصالًا --- وهو مخزن المعمَّر تحت الأرض.

\textbf{6.} بذورًا في التربة: وهي خطة الحولي --- \omterm{def:g1:plants-around-us:plant}{فالنباتات} نفسها
ماتت في الخريف.

\textbf{7.} السنونو: في إفريقيا، بذاته المهاجرة، متبعًا \omterm{prop:g3:sorting-living-things:groups}{الحشرات}
الطائرة. والنحل: في \omterm{def:g5:first-look-at-cells:cell}{الخلية}، عنقودًا شتويًّا، دافئًا حول ملكته.

\textbf{8.} البقاء نشطًا. وعلى العقعق أن يجد غذاءه عبر أشحّ
الشهور --- وهي المشكلة نفسها التي يفلت منها السنونو بالرحيل.

\textbf{9.} مثلًا: تنفتح \omterm{def:g6:life-through-seasons:buds}{البراعم} (آذار--نيسان) حين يبلغ النهار
$11$--$13$ ساعة وتتجاوز الصباحات
\qty{8}{\degreeCelsius}؛ وتذهّب \omterm{def:g1:plants-around-us:parts}{الأوراق} وتسقط (تشرين
الأول--تشرين الثاني) حين تقصر الأيام وتبرد الصباحات. فأحداث
الشجرة تقترن بتغيرات \omterm{def:g6:exploring-our-environment:conditions}{الشروط} صورةً بعد صورة.

\textbf{10.} الرقعة لم تمت: بل انسحبت إلى المخزن --- حية أبصالًا
تحت البقعة نفسها. و``ذهب'' تستحق دائمًا: ذهب \emph{إلى أين}،
وبأي \emph{صورة}؟

\textbf{11.} \omterm{def:g1:plants-around-us:tree}{أغصان} عارية تحمل \omterm{def:g6:life-through-seasons:buds}{براعم}؛ وعقعق في مكان ما على
الأرجح؛ وتحتها فرش \omterm{def:g1:plants-around-us:parts}{أوراق} وأرض نرجس عارية تخفي أبصالًا؛ \omterm{def:g2:from-seed-to-plant:seed}{وبذور} خشخاش
في تراب الممر. والرهانان الأكيدان هما \omterm{def:g6:life-through-seasons:buds}{البراعم} والعري --- فهما
يلزمان من آلة الشجرة \omterm{def:g6:life-through-seasons:trees}{النفضية} ومن برهان سنة كاملة؛ أما العقعق
فمرجّح لا مضمون.

\textbf{12.} يدل على أن الدفء كان الشرط الناقص --- \omterm{def:g6:exploring-our-environment:conditions}{فالضوء} داخل
البيت في شباط ليس أطول منه في الخارج، لكن الماء والغرفة دافئان،
وقد كفى ذلك. والاختبار المنصف: غصنان من الشجرة نفسها، كلاهما في
ماء عند النافذة نفسها، أحدهما في غرفة دافئة والآخر في رواق بارد
--- فرق واحد هو الدفء؛ فإن أورق الدافئ وحده ثبت الحكم ثبوتًا
قاطعًا.
\end{solution}
